\section{Methods}\label{section:methods}

\subsection{Initial definitions}

First, we introduce some basic components as part of QML. One-qubit bases $\ket{0}$ and $\ket{1}$ are defined in \cref{single-qubits}.

\begin{equation}\label{single-qubits}
    \ket{0} = \begin{pmatrix}
        1\\0
    \end{pmatrix}\quad
    \ket{1} = \begin{pmatrix}
        0\\1
    \end{pmatrix}
\end{equation}

Using these one-qubit bases, we can write quantum states as a linear combination as shown in \cref{quantum-states}, or using the Bloch-sphere representation rewritten as in \cref{bloch-sphere}
\begin{equation}\label{quantum-states}
    \ket{\psi} = \alpha\ket{0}+\beta\ket{1}
\end{equation}
\begin{equation}\label{bloch-sphere}
    \ket{\psi}=\cos\left(\frac{\theta}{2}\right)\ket{0}+\sin\left(\frac{\theta}{2}\right)e^{i\phi}\ket{1}
\end{equation}

Quantum algorithms are typically formulated using Pauli operators $\sigma_x,\space\sigma_y,\space\sigma_z$ as seen in \cref{pauli-matrices} and a small universal set of quantum gates that act on qubits. The Pauli matrices together with the identity $\textbf{I}$, form a convenient operator basis for expressing quantum Hamiltonians and observables. 
\begin{equation}\label{pauli-matrices}
    \sigma_x=\begin{pmatrix}
        0&1\\1&0
    \end{pmatrix}\quad
    \sigma_y=\begin{pmatrix}
        0&-i\\i&0
    \end{pmatrix}\quad
    \sigma_z=\begin{pmatrix}
        1&0\\0&-1
    \end{pmatrix}
\end{equation}

In practical implementations, these operators are realized through elementary gates such as the Hadamard gate ($H$), which creates superposition, the phase gate ($S$), which introduces controlled phase shifts, and the CNOT gate, which generates entanglement between qubits. All as seen in \cref{quantum-gates}. These gates form a standard set of building blocks for constructing quantum circuits and implementing algorithms such as the VQE.
\begin{equation}\label{quantum-gates}
\begin{array}{cc}
H = \frac{1}{\sqrt{2}}
\begin{pmatrix}
1 & 1\\
1 & -1
\end{pmatrix}
&
S =
\begin{pmatrix}
1 & 0\\
0 & -i
\end{pmatrix}
\\[12pt]
\multicolumn{2}{c}{
\text{CNOT} =
\begin{pmatrix}
1 & 0 & 0 & 0\\
0 & 1 & 0 & 0\\
0 & 0 & 0 & 1\\
0 & 0 & 1 & 0
\end{pmatrix}}
\end{array}
\end{equation}

\subsection{Tensor products and two-qubit states}
When describing systems of multiple qubits, the Hilbert space of the composite system is constructed using tensor (Kronecker) products of the individual qubit states as seen in \cref{kronecker-products}. This allows basis states of multi-qubit systems to be written as products of the single-qubit computational basis states $\ket{0}$ and $\ket{1}$. For a two-qubit system the computational basis therefore consists of the four states $\{\ket{00},\ket{01},\ket{10},\ket{11}\}$, obtained through tensor products of the single-qubit states.
\begin{equation}\label{kronecker-products}
\begin{aligned}
\ket{00} &= \ket{0} \otimes \ket{0} &
\ket{01} &= \ket{0} \otimes \ket{1} \\
\ket{10} &= \ket{1} \otimes \ket{0} &
\ket{11} &= \ket{1} \otimes \ket{1}
\end{aligned}
\end{equation}

Certain superpositions of these two-qubit basis states form maximally entangled states known as Bell states as seen in \cref{bell-states}. These states play an important role in quantum information theory and serve as fundamental examples of quantum entanglement between two qubits. 
\begin{equation}\label{bell-states}
\begin{aligned}
\ket{\Phi^+} &= \frac{\ket{00} + \ket{11}}{\sqrt{2}} &
\ket{\Phi^-} &= \frac{\ket{00} - \ket{11}}{\sqrt{2}} \\
\ket{\Psi^+} &= \frac{\ket{01} + \ket{10}}{\sqrt{2}} &
\ket{\Psi^-} &= \frac{\ket{01} - \ket{10}}{\sqrt{2}}
\end{aligned}
\end{equation}

For composite quantum systems, operators acting on multiple qubits are constructed using tensor products of single-qubit operators. If an operator $A$ acts on the first qubit and an operator $B$ acts on the second qubit, the combined operator acting on the two-qubit system is written as $A \otimes B$. In this way, Hamiltonians and quantum gates for multi-qubit systems can be expressed in terms of tensor products of Pauli matrices and the identity operator.

As an example, the operator $I \otimes \sigma_x$ flips the second qubit while leaving the first qubit unchanged. In the computational basis $\{\ket{00},\ket{01},\ket{10},\ket{11}\}$ this operator is given in \cref{I-otimes-sigma}.
\begin{equation}\label{I-otimes-sigma}
I \otimes \sigma_x =
\begin{pmatrix}
0 & 1 & 0 & 0 \\
1 & 0 & 0 & 0 \\
0 & 0 & 0 & 1 \\
0 & 0 & 1 & 0
\end{pmatrix}.
\end{equation}

Similarly, tensor products between Pauli operators act on both qubits simultaneously. For example, the operator $\sigma_x \otimes \sigma_x$ flips both qubits as seen in \cref{sigmax-otimes-sigmax}.
\begin{equation}\label{sigmax-otimes-sigmax}
\sigma_x \otimes \sigma_x =
\begin{pmatrix}
0 & 0 & 0 & 1 \\
0 & 0 & 1 & 0 \\
0 & 1 & 0 & 0 \\
1 & 0 & 0 & 0
\end{pmatrix}.
\end{equation}

While pure quantum states are written as state vectors, a more general description is given by density matrices. For a pure state $\ket{\psi}$ the density operator is given in \cref{density-matrix} and for a statistical ensemble $\{\ket{\psi_i},p_i\}$ it becomes as in \cref{mixed-density}.
\begin{equation}\label{density-matrix}
\rho = \ket{\psi}\bra{\psi},
\end{equation}
\begin{equation}\label{mixed-density}
\rho = \sum_i p_i \ket{\psi_i}\bra{\psi_i}.
\end{equation}

For bipartite systems the state of a subsystem is obtained via the partial trace, in  \cref{partial-trace}, which is central when analysing entanglement. A standard measure of quantum uncertainty is the von Neumann entropy in \cref{von-neumann-entropy}, which vanishes for pure states and is positive for mixed states.
\begin{equation}\label{partial-trace}
\rho_A = \mathrm{Tr}_B(\rho_{AB}),
\end{equation}
\begin{equation}\label{von-neumann-entropy}
\mathcal{S}(\rho) = -\mathrm{Tr}(\rho \log \rho),
\end{equation}

\subsection{Eigenvalue analysis}

The stationary Schrödinger equation defines the energy spectrum of a quantum system. The Hamiltonian operator acting on a quantum state produces the same state scaled by its corresponding energy eigenvalue, as seen in \cref{h-psi-e-psi}, where $\mathcal{H}$ is the Hamiltonian and $E$ the energy eigenvalue and $\ket{\psi}$ the corresponding eigenstate.
\begin{equation}\label{h-psi-e-psi}
    \mathcal{H}\ket{\psi}=E\ket{\psi}
\end{equation}

In the computational basis the Hamiltonian can be represented as a matrix acting on the state vector, and for a two-level system this can take the form in \cref{hamiltonian-vector}. Here the diagonal elements correspond to energies of basis states while the off-diagonal elements represents interactions between them.
\begin{equation}\label{hamiltonian-vector}
    \mathcal{H}=\begin{bmatrix}
        \mathcal{H}_{11} & \mathcal{H}_{12} \\
        \mathcal{H}_{21} & \mathcal{H}_{22}
    \end{bmatrix}
\end{equation}

Furthermore, the Hamiltonian is decomposed into $\mathcal{H}=\mathcal{H}_0+\mathcal{H}_1$, with a solvable non-interacting component $\mathcal{H}_0$ and an interaction term $\mathcal{H}_1$. See \cref{H0-H1}, where $E_i$ represent the unperturbed energy levels and $V_{ij}$ describe interaction matrix elements.
\begin{equation}\label{H0-H1}
    \mathcal{H}_0=\begin{bmatrix}
        E_1 & 0 \\
        0 & E_2
    \end{bmatrix}\quad
    \mathcal{H}_1=\begin{bmatrix}
        V_{11} & V_{12} \\
        V_{21} & V_{22}
    \end{bmatrix}
\end{equation} 

Any $2\times2$ Hamiltonian can be expressed using the Pauli operator basis together with the identity operator. For the present system this leads to the identities seen in \cref{H0-H1-decomp} with \cref{mathcal-E-Omega}. 
\begin{equation}\label{H0-H1-decomp}
\begin{aligned}
    &\mathcal{H}=a\textbf{I}+b\sigma_x+c\sigma_y+d\sigma_z\\
    &\mathcal{H}_0 = \mathcal{E}\textbf{I} + \Omega\sigma_z \\
    &\mathcal{H}_1 = c\textbf{I} + \omega_z\sigma_z + \omega_x\sigma_x
\end{aligned}
\end{equation}
\begin{equation}\label{mathcal-E-Omega}
    \mathcal{E}=\frac{E_1+E_2}{2} \quad \Omega=\frac{E_1-E_2}{2}
\end{equation}

Now, we can move on to the analytical eigenvalue solution, where the eigenvalues of the Hamiltonian are obtained from the characteristic equation in \cref{characteristic-equation}.
\begin{equation}\label{characteristic-equation}
    \det(\mathcal{H}-E\textbf{I})=0
\end{equation}

For a general $2\times2$ Hamiltonian in \cref{general-hamiltonian} the resulting energy eigenvalues can be found in \cref{energy-eigenvalues}. Using this, the lowest eigenvalue corresponds to the ground state energy of the system.
\begin{equation}\label{general-hamiltonian}
    \mathcal{H}=\begin{pmatrix}
        a&b\\b&d
    \end{pmatrix}
\end{equation}
\begin{equation}\label{energy-eigenvalues}
    E_{\pm} = \frac{a+d}{2}\pm\sqrt{\left(\frac{a-d}{2}\right)^2+b^2}
\end{equation} 

To study the effects of interactions in \cref{hamilton-interaction}, the Hamiltonian is parametrised by a coupling parameter $\lambda\in[0,1]$, where the limits $\lambda=0$ and $\lambda=1$ correspond to the non-interacting and fully interacting systems respectively.
\begin{equation}\label{hamilton-interaction}
    \mathcal{H}=\mathcal{H}_0+\lambda \mathcal{H}_1
\end{equation}

The Hamiltonian is diagonalised numerically to obtain its eigenvalues and eigenvectors. For Hermitian matrices this is efficiently performed using the \texttt{numpy} package in \texttt{python}, as seen in \cref{lst:numpy-diagonalization}.\newline

\begin{lstlisting}[
caption={Diagonalization of the Hamiltonian matrix $\mathcal{H}$ using NumPy's \texttt{np.linalg.eigh}, which computes the eigenvalues and eigenvectors of a Hermitian operator. This routine is used as the classical baseline for validating the quantum variational results.},
label={lst:numpy-diagonalization},
captionpos=b
]
np.linalg.eigh(H)
\end{lstlisting}

\subsection{Variational Quantum Eigensolver (VQE)}

VQE is based on the variational principle, for any normalised trial state $\ket{\psi(\theta)}$, the expectation value of the Hamiltonian provides an upper bound to the true ground state energy. The goal in using a VQE is therefore to minimize the energy expectation value ewith respect to the circuit parameter $\theta$ shown in \cref{minimum-expectation}. 
\begin{equation}\label{minimum-expectation}
    E_0\leq\min_\theta\bra{\psi(\theta)}\mathcal{H}\ket{\psi(\theta)}
\end{equation}

A parameterized quantum circuit $U(\theta)$ prepares a trial state from a simple initial computational basis state. By varying the parameters $\theta$, the circuit explores different candidate states in \cref{parameterized-psi}
\begin{equation}\label{parameterized-psi}
    \ket{\psi(\theta)}=U(\theta)\ket{0}^{\otimes n}
\end{equation}

To evaluate the energy on a quantum computer, the Hamiltonian is expressed as a weighted sum of tensor products of Pauli operators, in \cref{c-P-pauli}. Each term can then be measured seperately on the quantum circuit, where $c_k$ are the real coefficients and $P_k$ are Pauli strings acting on the qubits.
\begin{equation}\label{c-P-pauli}
    \mathcal{H}=\sum_k c_kP_k 
    \quad
    P_k\in\{I,X,Y,Z\}^{\otimes n}
\end{equation}

The VQE cost function is the expectation value of the Hamiltonian with respect to the trial state. Using the Pauli decomposition, the energy can be written as a sum of expectation values of Pauli operators in \cref{E-theta-bra-psi}.
\begin{equation}\label{E-theta-bra-psi}
    E(\theta)=\bra{\psi(\theta)}\mathcal{H}\ket{\psi(\theta)} = \sum_k c_k \langle P_k \rangle_\theta
\end{equation}

A classical optimizer updates the circuit parameter iteratively in order to reduce the energy expectation value. Conceptually, this can be written as an update rule, where $\eta$ is a learning rate or step size, as in \cref{nabla-theta}.
\begin{equation}\label{nabla-theta}
    \theta_{t+1} = \theta_t - \eta\nabla_{\theta}E(\theta_t)
\end{equation}

After convergence, the optimal parameters $\theta^*$ define the approximate ground state prepared by the circuit. The estimated ground state energy is then shown in \cref{ground-staet-E}.
\begin{equation}\label{ground-staet-E}
    E_{\text{VQE}} = E(\theta)
\end{equation}

The choice of parameterized circuit $U(\theta)$ determines the family of states explored by the VQE and therefore strongly influences the quality of the variational solution. For a single qubit, a minimal ansatz consists of a single rotation gate is seen in \cref{ry-ansatz}.
\begin{equation}\label{ry-ansatz}
\ket{\psi(\theta)} = R_y(\theta)\ket{0}
= \cos\frac{\theta}{2}\ket{0} + \sin\frac{\theta}{2}\ket{1}.
\end{equation}

For two qubits, a common hardware-efficient architecture combines single-qubit rotations with an entangling gate. A simple example is shown in \cref{two-qubit-ansatz}.
\begin{equation}\label{two-qubit-ansatz}
\ket{\psi(\boldsymbol{\theta})} =
\big(R_y(\theta_3)\otimes R_y(\theta_4)\big)
\,\mathrm{CNOT}\,
\big(R_y(\theta_1)\otimes R_y(\theta_2)\big)\ket{00}.
\end{equation}

\subsection{The Lipkin model}

The Lipkin--Meshkov--Glick (LMG) model describes $N$ fermions distributed over two $N$-fold degenerate energy levels separated by an energy $\varepsilon$. In terms of quasispin operators the Hamiltonian is written as in \cref{lipkin-hamiltonian}.
\begin{equation}\label{lipkin-hamiltonian}
    \mathcal{H} = \varepsilon J_z + \frac{1}{2}V(J_+^2 + J_-^2)
\end{equation}

The operators $J_z, J_+, J_-$ satisfy the $\mathfrak{su}(2)$ algebra and act within the angular momentum sector $J=N/2$. The first term represents the single-particle energy splitting, while the interaction term scatters particle pairs between the two levels. To simulate the model on a quantum computer, the quasispin operators are mapped to Pauli operators acting on $N$ qubits. This yields the qubit Hamiltonian shown in \cref{lipkin-pauli-hamiltonian}.
\begin{equation}\label{lipkin-pauli-hamiltonian}
    \mathcal{H} =
    \frac{\varepsilon}{2}\sum_{i=1}^{N} Z_i
    +
    \frac{V}{2}\sum_{i<j}(X_i X_j - Y_i Y_j)
\end{equation}

This Pauli representation allows the Lipkin Hamiltonian to be directly evaluated within the VQE framework introduced in the following subsection.

\subsection{Implementation in \texttt{python}}

All simulations and numerical verification routines were implemented in \texttt{python}. Linear algebra and numerical routines were handled using \texttt{numpy} and \texttt{scipy}, data handling with \texttt{pandas}, and visualization with \texttt{matplotlib} and \texttt{seaborn}. Quantum circuits and hybrid quantum–classical optimization were implemented using \texttt{PennyLane} and \texttt{Qiskit} \parencite{numpy,pennylane,qiskit}. Python has become a standard platform for numerical computing and quantum algorithm prototyping due to its extensive scientific ecosystem. 

The full implementation and supporting code are available in the project repository \parencite{fys5419-repo}, with the development environment summarized in \cref{lst:python-environment} in the Appendix. Parts of the implementation follow examples and structures from the quantum machine learning teaching repository by \textcite{compphysics_qcml_repo}.

\subsection{Use of AI-assisted development tools}

AI-assisted programming tools were used during development for code generation, debugging, and documentation. These included large language models such as \texttt{ChatGPT} \parencite{chatgpt}, \texttt{Claude} \parencite{claude}, \texttt{Grok} \parencite{grok}, and \texttt{Gemini} \parencite{gemini}, as well as coding assistants such as \texttt{GitHub Copilot}, \texttt{Antigravity} and \texttt{Cursor} \parencite{copilot,google-antigravity,cursor}. 
